\documentclass{amsart}
% For dealing with references we use the comment environment
\usepackage{verbatim}
\newenvironment{reference}{\comment}{\endcomment}
%\newenvironment{reference}{}{}
\newenvironment{slogan}{\comment}{\endcomment}
\newenvironment{history}{\comment}{\endcomment}

% For commutative diagrams we use Xy-pic
\usepackage[all]{xy}

% We use 2cell for 2-commutative diagrams.
\xyoption{2cell}
\UseAllTwocells

% We use multicol for the list of chapters between chapters
\usepackage{multicol}

% This is generally recommended for better output
\usepackage{lmodern}
\usepackage[T1]{fontenc}

% For cross-file-references

% Package for hypertext links:
\usepackage{hyperref}

% For any local file, say "hello.tex" you want to link to please

% Theorem environments.
%
\theoremstyle{plain}
\newtheorem{theorem}[subsection]{Theorem}
\newtheorem{proposition}[subsection]{Proposition}
\newtheorem{lemma}[subsection]{Lemma}

\theoremstyle{definition}
\newtheorem{definition}[subsection]{Definition}
\newtheorem{example}[subsection]{Example}
\newtheorem{exercise}[subsection]{Exercise}
\newtheorem{situation}[subsection]{Situation}

\theoremstyle{remark}
\newtheorem{remark}[subsection]{Remark}
\newtheorem{remarks}[subsection]{Remarks}

\numberwithin{equation}{subsection}

% Macros
%
\def\lim{\mathop{\mathrm{lim}}\nolimits}
\def\colim{\mathop{\mathrm{colim}}\nolimits}
\def\Spec{\mathop{\mathrm{Spec}}}
\def\Hom{\mathop{\mathrm{Hom}}\nolimits}
\def\Ext{\mathop{\mathrm{Ext}}\nolimits}
\def\SheafHom{\mathop{\mathcal{H}\!\mathit{om}}\nolimits}
\def\SheafExt{\mathop{\mathcal{E}\!\mathit{xt}}\nolimits}
\def\Sch{\mathit{Sch}}
\def\Mor{\mathop{\mathrm{Mor}}\nolimits}
\def\Ob{\mathop{\mathrm{Ob}}\nolimits}
\def\Sh{\mathop{\mathit{Sh}}\nolimits}
\def\NL{\mathop{N\!L}\nolimits}
\def\CH{\mathop{\mathrm{CH}}\nolimits}
\def\proetale{{pro\text{-}\acute{e}tale}}
\def\etale{{\acute{e}tale}}
\def\QCoh{\mathit{QCoh}}
\def\Ker{\mathop{\mathrm{Ker}}}
\def\Im{\mathop{\mathrm{Im}}}
\def\Coker{\mathop{\mathrm{Coker}}}
\def\Coim{\mathop{\mathrm{Coim}}}

% Boxtimes
%
\DeclareMathSymbol{\boxtimes}{\mathbin}{AMSa}{"02}

%
% Macros for moduli stacks/spaces
%
\def\QCohstack{\mathcal{QC}\!\mathit{oh}}
\def\Cohstack{\mathcal{C}\!\mathit{oh}}
\def\Spacesstack{\mathcal{S}\!\mathit{paces}}
\def\Quotfunctor{\mathrm{Quot}}
\def\Hilbfunctor{\mathrm{Hilb}}
\def\Curvesstack{\mathcal{C}\!\mathit{urves}}
\def\Polarizedstack{\mathcal{P}\!\mathit{olarized}}
\def\Complexesstack{\mathcal{C}\!\mathit{omplexes}}
% \Pic is the operator that assigns to X its picard group, usage \Pic(X)
% \Picardstack_{X/B} denotes the Picard stack of X over B
% \Picardfunctor_{X/B} denotes the Picard functor of X over B
\def\Pic{\mathop{\mathrm{Pic}}\nolimits}
\def\Picardstack{\mathcal{P}\!\mathit{ic}}
\def\Picardfunctor{\mathrm{Pic}}
\def\Deformationcategory{\mathcal{D}\!\mathit{ef}}

\setlength{\emergencystretch}{3em}
\begin{document}
\title[Stacks Project, Tag 07AC]{660 --- Derived lower shriek exists for continuous cocontinuous site functors}
\maketitle
\noindent Copyright (C) 2005--2025 Johan de Jong. Stacks Project authors. Source: \href{https://stacks.math.columbia.edu/tag/07AC}{Tag 07AC}.

\noindent Editorial difficulty note (not author text). Difficulty H2: The proof overcomes failure of pullback to preserve injectives using acyclicity on extension-by-zero generators and an Ext spectral sequence. It constructs an unbounded left-derived adjoint through a genuinely nonexact functor..

\noindent Editorial provenance note (not author text). History: excerpt from revision \texttt{a0444\allowbreak 6e57e\allowbreak c1fbc\allowbreak 252a8\allowbreak 71afc\allowbreak ec775\allowbreak 2fb28\allowbreak 07b14}, sites-cohomology.tex, lines 9309--9433. Reference presentation was adapted; original-author context and core support blocks were retained or added. The mathematical wording is unchanged. Licensed under GNU FDL 1.2 or later; no invariant sections or cover texts. See ../licenses/COPYING. Context blocks reproduce author definitions and supporting results; they are not additional counted problems.

\begin{lemma}
\label{lemma-pullback-injective-pre-limp}
Let $u : \mathcal{C} \to \mathcal{D}$ be a continuous and cocontinuous
functor of sites. Let $g : \Sh(\mathcal{C}) \to \Sh(\mathcal{D})$
be the corresponding morphism of topoi. Let $\mathcal{O}_\mathcal{D}$
be a sheaf of rings and let $\mathcal{I}$ be an injective
$\mathcal{O}_\mathcal{D}$-module. Then
$H^p(U, g^{-1}\mathcal{I}) = 0$ for all $p > 0$ and $U \in \Ob(\mathcal{C})$.
\end{lemma}

\begin{proof}
The vanishing of the lemma follows from
Lemma \href{https://stacks.math.columbia.edu/tag/03F9}{lemma cech vanish collection (Tag 03F9)}
if we can prove vanishing of all higher
{\v C}ech cohomology groups
$\check H^p(\mathcal{U}, g^{-1}\mathcal{I})$
for any covering $\mathcal{U} = \{U_i \to U\}$ of $\mathcal{C}$.
Since $u$ is continuous, $u(\mathcal{U}) = \{u(U_i) \to u(U)\}$
is a covering of $\mathcal{D}$, and
$u(U_{i_0} \times_U \ldots \times_U U_{i_n}) =
u(U_{i_0}) \times_{u(U)} \ldots \times_{u(U)} u(U_{i_n})$.
Thus we have
$$
\check H^p(\mathcal{U}, g^{-1}\mathcal{I}) =
\check H^p(u(\mathcal{U}), \mathcal{I})
$$
because $g^{-1} = u^p$ by Sites, Lemma \href{https://stacks.math.columbia.edu/tag/00XR}{lemma when shriek (Tag 00XR)}.
Since $\mathcal{I}$ is an injective
$\mathcal{O}_\mathcal{D}$-module these {\v C}ech cohomology groups vanish, see
Lemma \href{https://stacks.math.columbia.edu/tag/03FC}{lemma injective module trivial cech (Tag 03FC)}.
\end{proof}


\begin{lemma}
\label{lemma-existence-derived-lower-shriek}
Let $u : \mathcal{C} \to \mathcal{D}$ be a continuous and cocontinuous
functor of sites. Let $g : \Sh(\mathcal{C}) \to \Sh(\mathcal{D})$ be the
corresponding morphism of topoi. Let $\mathcal{O}_\mathcal{D}$
be a sheaf of rings and set
$\mathcal{O}_\mathcal{C} = g^{-1}\mathcal{O}_\mathcal{D}$.
The functor $g_! : \textit{Mod}(\mathcal{O}_\mathcal{C}) \to
\textit{Mod}(\mathcal{O}_\mathcal{D})$
(see
Modules on Sites, Lemma \href{https://stacks.math.columbia.edu/tag/0797}{lemma lower shriek modules (Tag 0797)})
has a left derived functor
$$
Lg_! : D(\mathcal{O}_\mathcal{C}) \longrightarrow D(\mathcal{O}_\mathcal{D})
$$
which is left adjoint to $g^*$. Moreover, for $U \in \Ob(\mathcal{C})$ we
have
$$
Lg_!(j_{U!}\mathcal{O}_U) =
g_!j_{U!}\mathcal{O}_U =
j_{u(U)!} \mathcal{O}_{u(U)}.
$$
where $j_{U!}$ and $j_{u(U)!}$ are extension by zero associated to the
localization morphism
$j_U : \mathcal{C}/U \to \mathcal{C}$ and
$j_{u(U)} : \mathcal{D}/u(U) \to \mathcal{D}$.
\end{lemma}

\begin{proof}
We are going to use
Derived Categories, Proposition \href{https://stacks.math.columbia.edu/tag/0794}{proposition left derived exists (Tag 0794)}
to construct $Lg_!$. To do this we have to verify assumptions
(1), (2), (3), (4), and (5) of that proposition.
First, since $g_!$ is a left adjoint
we see that it is right exact and commutes with all colimits, so
(5) holds. Conditions (3) and (4) hold because the category of modules
on a ringed site is a Grothendieck abelian category.
Let $\mathcal{P} \subset \Ob(\textit{Mod}(\mathcal{O}_\mathcal{C}))$
be the collection of $\mathcal{O}_\mathcal{C}$-modules which are direct
sums of modules of the form $j_{U!}\mathcal{O}_U$. Note that
$g_!j_{U!}\mathcal{O}_U = j_{u(U)!} \mathcal{O}_{u(U)}$, see proof of
Modules on Sites, Lemma \href{https://stacks.math.columbia.edu/tag/0797}{lemma lower shriek modules (Tag 0797)}.
Every $\mathcal{O}_\mathcal{C}$-module is a quotient of an object of
$\mathcal{P}$, see
Modules on Sites, Lemma \href{https://stacks.math.columbia.edu/tag/03EW}{lemma module quotient flat (Tag 03EW)}.
Thus (1) holds. Finally, we have to prove (2).
Let $\mathcal{K}^\bullet$ be a bounded above acyclic complex of
$\mathcal{O}_\mathcal{C}$-modules with $\mathcal{K}^n \in \mathcal{P}$
for all $n$. We have to show that $g_!\mathcal{K}^\bullet$ is
exact. To do this it suffices to show, for every injective
$\mathcal{O}_\mathcal{D}$-module $\mathcal{I}$ that
$$
\Hom_{D(\mathcal{O}_\mathcal{D})}(
g_!\mathcal{K}^\bullet, \mathcal{I}[n]) = 0
$$
for all $n \in \mathbf{Z}$. Since $\mathcal{I}$ is injective we have
\begin{align*}
\Hom_{D(\mathcal{O}_\mathcal{D})}(
g_!\mathcal{K}^\bullet, \mathcal{I}[n])
& =
\Hom_{K(\mathcal{O}_\mathcal{D})}(
g_!\mathcal{K}^\bullet, \mathcal{I}[n]) \\
& =
H^n(\Hom_{\mathcal{O}_\mathcal{D}}(
g_!\mathcal{K}^\bullet, \mathcal{I})) \\
& =
H^n(\Hom_{\mathcal{O}_\mathcal{C}}(
\mathcal{K}^\bullet, g^{-1}\mathcal{I}))
\end{align*}
the last equality by the adjointness of $g_!$ and $g^{-1}$.

\medskip\noindent
The vanishing of this group would be clear if $g^{-1}\mathcal{I}$
were an injective $\mathcal{O}_\mathcal{C}$-module. But
$g^{-1}\mathcal{I}$ isn't necessarily an injective
$\mathcal{O}_\mathcal{C}$-module as $g_!$ isn't exact in
general. We do know that
$$
\Ext^p_{\mathcal{O}_\mathcal{C}}(
j_{U!}\mathcal{O}_U, g^{-1}\mathcal{I}) =
H^p(U, g^{-1}\mathcal{I}) = 0 \text{ for }p \geq 1
$$
Here the first equality follows from
$\Hom_{\mathcal{O}_\mathcal{C}}(j_{U!}\mathcal{O}_U, \mathcal{H}) =
\mathcal{H}(U)$ and taking derived functors and the vanishing of
$H^p(U, g^{-1}\mathcal{I})$ for $p > 0$ and $U \in \Ob(\mathcal{C})$
follows from Lemma \ref{lemma-pullback-injective-pre-limp}.
Since each $\mathcal{K}^{-q}$ is a direct sum of modules of the form
$j_{U!}\mathcal{O}_U$ we see that
$$
\Ext^p_{\mathcal{O}_\mathcal{C}}(\mathcal{K}^{-q}, g^{-1}\mathcal{I}) = 0
\text{ for }p \geq 1\text{ and all }q
$$
Let us use the spectral sequence (see
Example
\href{https://stacks.math.columbia.edu/tag/07AA}{example hom complex into sheaf (Tag 07AA)})
$$
E_1^{p, q} = \Ext^q_{\mathcal{O}_\mathcal{C}}(
\mathcal{K}^{-p}, g^{-1}\mathcal{I})
\Rightarrow
\Ext^{p + q}_{\mathcal{O}_\mathcal{C}}(
\mathcal{K}^\bullet, g^{-1}\mathcal{I}) = 0.
$$
Note that the spectral sequence abuts to zero as $\mathcal{K}^\bullet$
is acyclic (hence vanishes in the derived category, hence produces
vanishing ext groups). By the vanishing of higher exts proved above
the only nonzero terms on the $E_1$ page are the terms
$E_1^{p, 0} = \Hom_{\mathcal{O}_\mathcal{C}}(
\mathcal{K}^{-p}, g^{-1}\mathcal{I})$.
We conclude that the complex
$\Hom_{\mathcal{O}_\mathcal{C}}(
\mathcal{K}^\bullet, g^{-1}\mathcal{I})$
is acyclic as desired.

\medskip\noindent
Thus the left derived functor $Lg_!$ exists.
It is left adjoint to $g^{-1} = g^* = Rg^* = Lg^*$, i.e., we have
\begin{equation}
\label{equation-to-prove}
\Hom_{D(\mathcal{O}_\mathcal{C})}(K, g^*L) =
\Hom_{D(\mathcal{O}_\mathcal{D})}(Lg_!K, L)
\end{equation}
by Derived Categories, Lemma \href{https://stacks.math.columbia.edu/tag/09T5}{lemma derived adjoint functors (Tag 09T5)}.
This finishes the proof.
\end{proof}
\end{document}
